\emph{(i).} With $E = \hb^{-1}(a)$, a spacetime vector, and $u\cdot E = 0$, $u\wedge E = uE$. For a vector $W$ of the added plane, $(uE)\times W = \tfrac12(uEW - WuE) = -(u\cdot W)E$, since $E$ anticommutes with $u$ and $W$. With $u\cdot\nInf = -\rho$ and $u\cdot\nO = 1$ this gives $-(u\wedge E)\times(\sigma\nInf + \tau\nO) = (\tau - \rho\sigma)E$, and $\Omega(a)$, a spacetime bivector, commutes with the added plane. For a spacetime vector $\mu$, $(uE)\times\mu = \tfrac12 u(E\mu + \mu E) = (E\cdot\mu)u$. In particular $A_{\rho}(a)\times I_{\rho} = (\rho - \rho)E = 0$, so $D_{a}I_{\rho} = 0$. The field strength is $F(a\wedge b) = a\cdot\nabla A_{\rho}(b) - b\cdot\nabla A_{\rho}(a) + A_{\rho}(a)\times A_{\rho}(b)$. Its part in $u\wedge(\text{spacetime})$ is $-u\wedge$ (the torsion), and $(u\wedge E_{a})\times(u\wedge E_{b}) = -Q(u)\,E_{a}\wedge E_{b} = 2\rho\,E_{a}\wedge E_{b}$, so for vanishing torsion $F$ is the field strength of $\Omega$ plus $2\rho\,\hb^{-1}(a)\wedge\hb^{-1}(b)$, which on $a\wedge b = \hb(B)$ is $-(\mathcal{R}(B) - 2\rho B)$ in the store's sign; this is clause (ii) of \claimref{synthesis.frozen-point-field-space-is-the-gtg-pair-and-its-omega-equation-is-the-torsion-equation}. A spacetime bivector commutes with $\nInf$ and $\nO$, so $F\times V = L\cdot\mu$. The check computes $F$ from the fields themselves on de Sitter, six isotropic fluid instants and a Schwarzschild field in Kerr-Schild form and compares.

\emph{(ii).} If $D_{a}V = 0$ for all $a$ on $U$, then $F(a\wedge b)\times V = [D_{a}, D_{b}]V - D_{[a,b]}V = 0$, so $L(B)\cdot\mu = 0$ for all $B$ at every point. The field strength is continuous, and having no kernel vector is an open condition (the rank of a continuous family of linear maps is lower semicontinuous), so it holds on a neighbourhood $U_{0}$ of the point, and $\mu = 0$ on $U_{0}$. There, by (i), $D_{a}V = (a\cdot\nabla\sigma)\nInf + (a\cdot\nabla\tau)\nO + (\tau - \rho\sigma)\hb^{-1}(a)$. The three parts are linearly independent and $\hb^{-1}(a) \neq 0$ for $a \neq 0$ on the invertible region, so $\sigma$ and $\tau$ are constant on $U_{0}$ and $\tau = \rho\sigma$: $V = \sigma I_{\rho}$ on $U_{0}$. Then $V - \sigma I_{\rho}$ is parallel on $U$ and vanishes on $U_{0}$; a parallel section of a linear connection that vanishes at a point vanishes along every path from it, and $U$ is connected, so $V = \sigma I_{\rho}$ on $U$. Nothing used $\rho \neq 0$.

\emph{(iii).} On an Einstein field with $\mathcal{R}(a) = \Lambda a$ the curvature operator is $\mathcal{W} + \tfrac{\Lambda}{3}\mathrm{id}$ (its Ricci part is that of the de Sitter field), so $L = -\mathcal{W}$, and a kernel vector of $L$ is a kernel vector of the Weyl part, the condition of \claimref{synthesis.frozen-point-unique-where-weyl-has-no-kernel}; the check confirms $F = -\mathcal{W}$ on the Schwarzschild field at $\rho = 0$. On the de Sitter field at that $\Lambda$, $\mathcal{R}(B) = \tfrac{\Lambda}{3}B$ and $F = 0$: the connection is flat and parallel transport from a point gives a parallel tractor for each of the six initial values on a simply connected neighbourhood.

\emph{(iv).} On the isotropic field with flat slices the curvature operator is diagonal on the frame bivectors, with one value on the three time-space bivectors, fixed by $\ddot a/a$, and one on the three space-space bivectors, fixed by $H^{2}$; on the de Sitter field both are fixed by $\Lambda/3$. Subtracting, the two values of $L$ are proportional to $\ddot a/a - \Lambda/3 = -\tfrac16\kappa\rho_{m}(1 + 3w)$ and to $H^{2} - \Lambda/3 = \tfrac13\kappa\rho_{m}$. If both are nonzero, $L(\gamma_{0}\wedge\gamma_{i})\cdot\mu = 0$ and $L(\gamma_{i}\wedge\gamma_{j})\cdot\mu = 0$ force every component of $\mu$ to vanish. If the first vanishes and the second does not, the space-space bivectors kill exactly the multiples of $\gamma_{0}$. The check builds the instants from the two equations and finds kernel dimension $0$ at $w = 0, \tfrac13, 1, -1, -\tfrac23$ and $1$ at $w = -\tfrac13$.

\emph{(v).} With torsion, by the computation of (i), $F = L - u\wedge T$, and by the two actions of (i) $F\times V = L\cdot\mu - (T\cdot\mu)u + (\tau - \rho\sigma)T$. Given $\mu_{0} \neq 0$ and $c \neq 0$, the choice $T(B) = -(L(B)\cdot\mu_{0})/c$ makes the spacetime part vanish, and $T(B)\cdot\mu_{0} = -(L(B)\cdot\mu_{0})\cdot\mu_{0}/c = 0$ because $L(B)$ is a bivector; the check does this on a dust instant. If the stated map is injective at a point it is injective nearby, so $\mu = 0$ and $\tau = \rho\sigma$ there, and the derivative step of (ii) applies with $\tau - \rho\sigma = 0$ (argued; no field with that torsion is exhibited).

\emph{Applied to towers.} Immediate from (ii): $\ip{X}{cI_{\rho}} = c$ and $Q(cI_{\rho}) = 2\rho c^{2}$. \emph{Remark.} In the reversed form $x\times' y = -(x\times y)$ for a bivector $x$ and a vector $y$, and the pairing of $\nO$ with $\nInf$ changes sign; the two renormalisations and the conjugation by $I_{4}$ are computed in the check.
